% ===
\section{Research Design}
% ===

Building on the conceptual discussion presented in the Introduction, this section outlines the research design of the thesis. The study combines theoretical analysis, formal modeling, and empirical investigation to examine whether environmental preservation can be reframed not merely as an external restriction on production, but as an endogenous and incentive-compatible component of economic decision-making.

The research is guided by the concept of \emph{economic commensalism}, understood here as a theoretical framework in which ecological preservation may become compatible with private productivity and profitability when incentives are spatially grounded, intertemporally coherent, and strategically stable. The central purpose is not to reject existing environmental instruments, such as regulation, enforcement, carbon markets, payments for environmental services, or community governance, but to investigate how these instruments can be reorganized around the structural payoff of economic agents.

The thesis is organized into three complementary chapters. Chapter 1 develops the theoretical foundations of economic commensalism, discussing the relationship between production theory, environmental preservation, extractivism, predation, common-pool resources, and incentive-compatible governance. Chapter 2 formalizes the proposed framework through a dynamic model that incorporates environmental variables, intertemporal choice, uncertainty, bounded rationality, learning, and strategic interaction. Chapter 3 evaluates the empirical feasibility of the framework through an application to the Brazilian state of Mato Grosso, a territory marked by agricultural expansion, deforestation pressures, and ecological heterogeneity across biomes.

Together, these three components seek to provide conceptual clarity, formal consistency, and empirical plausibility to the central claim of the thesis: environmental preservation may become economically rational when ecological variables are incorporated into production decisions and when institutional incentives are designed according to spatial, temporal, and behavioral constraints.

% ---
\subsection{Research Problem}
% ---

The research problem emerges from a specific limitation in conventional approaches to environmental economics and policy design. Although environmental degradation is widely recognized as a problem of externalities, common-pool resources, institutional failure, and intertemporal welfare, environmental preservation is still frequently treated as an external obligation imposed on private agents. In this setting, producers often face preservation as a legal constraint, compliance cost, reputational requirement, or compensatory burden, rather than as a component of their own productive strategy.

This thesis therefore asks whether the apparent conflict between private profitability and socially desirable environmental preservation is inevitable, or whether it results, at least in part, from theoretical and institutional misalignment. If ecological variables are treated as exogenous to production, and if environmental policy operates mainly as a corrective layer imposed after private decisions have already been made, preservation will tend to appear as a burden. Conversely, if ecological preservation is incorporated into the productive calculus of agents, it may become possible to design institutional arrangements in which preserving environmental assets is compatible with productivity, profitability, and long-term welfare.

Accordingly, the central research question is:

\vspace{1em}

\begin{flushright}
\begin{minipage}{0.93\linewidth}
\noindent
\hspace{0.4em}\vrule width 0.4pt height 5.5em \hspace{1.2em}%
\parbox[b]{0.95\linewidth}{
\textit{Can environmental preservation be theoretically reframed as an endogenous and incentive-compatible component of economic decision-making, capable of reconciling private productivity and profitability with socially desirable ecological outcomes?}
}
\end{minipage}
\end{flushright}

\vspace{1em}

This question guides the theoretical, formal, and empirical investigation developed throughout the thesis. An affirmative answer would indicate that environmental preservation can be designed as part of the agent's rational economic calculus, rather than depending exclusively on coercion, moral commitment, or fragmented compensation. A negative or partial answer would still be relevant, since it would clarify the limits of incentive-compatible environmental governance and identify the conditions under which preservation cannot be reconciled with private profitability without stronger forms of regulation or redistribution.

% ---
\subsection{Hypotheses}
% ---

The hypotheses of the thesis are structured around one general hypothesis and five specific hypotheses. The general hypothesis expresses the central proposition of the research. The specific hypotheses decompose this proposition into theoretical, behavioral, spatial, institutional, and strategic dimensions.

%...
\subsubsection{General Hypothesis}
%...

The general hypothesis is that:

\begin{itemize}

    \item[\textbf{H0.}]
    
    \textit{It is possible to design an economic-institutional framework in which environmental preservation becomes compatible with private productivity and profitability when ecological variables are endogenized into production decisions and incentives are spatially, intertemporally, and strategically calibrated.}

\end{itemize}

H0 does not assume that preservation and profitability are automatically compatible. Rather, it states that compatibility may emerge under specific theoretical and institutional conditions. These conditions are examined through the specific hypotheses below.

%...
\subsubsection{Specific Hypotheses}
%...

\begin{itemize}

    \item[\textbf{H1.}]
    
    \textit{Endogenizing environmental variables into the production function is necessary to correct the underrepresentation of ecological systems in production theory and to improve the design of environmental governance mechanisms.}

\end{itemize}

H1 addresses the theoretical foundation of the thesis. If the environment is treated only as an externality, constraint, or fixed stock, its effects on productivity, resilience, and intertemporal welfare tend to be represented only after degradation occurs. By incorporating ecological variables into production decisions, the thesis investigates whether environmental preservation can be modeled as part of the productive structure itself.

\begin{itemize}

    \item[\textbf{H2.}]
    
    \textit{Economic decisions regarding environmental exploitation or preservation are intertemporal and shaped by bounded rationality, incomplete information, uncertainty, and adaptive learning.}

\end{itemize}

H2 defines the behavioral and temporal nature of the problem. Environmental decisions are not purely static choices made under perfect information. Agents face uncertain ecological feedback, price fluctuations, institutional instability, credit constraints, enforcement uncertainty, and imperfect knowledge about long-term consequences. This makes the preservation--exploitation dilemma dynamic and requires a framework capable of representing learning, expectations, and intertemporal trade-offs.

\begin{itemize}

    \item[\textbf{H3.}]
    
    \textit{Environmental incentives are more likely to be effective when they are spatially and ecologically targeted to the units affected by production, such as biomes, watersheds, ecoregions, or transition zones.}

\end{itemize}

H3 addresses the spatial dimension of environmental governance. Ecological degradation is territorially specific, while many environmental compensation mechanisms are standardized, centralized, or globally tradable. The hypothesis is that incentives detached from the affected ecological unit may satisfy accounting or regulatory requirements without preserving the ecological functions that sustain local production and resilience.

\begin{itemize}

    \item[\textbf{H4.}]
    
    \textit{Properly calibrated positive or mixed incentives can alter the structural payoff of economic agents, making preservation more attractive than degradation under specific institutional and economic conditions.}

\end{itemize}

H4 addresses the institutional mechanism of the thesis. The argument is not that positive incentives are always superior to negative incentives, nor that coercion is unnecessary. Rather, the hypothesis is that incentives linked to credit, certification, fiscal treatment, market access, carbon valuation, production rights, or payments for environmental services may reduce dependence on punitive enforcement when they make preservation privately advantageous.

\begin{itemize}

    \item[\textbf{H5.}]
    
    \textit{Under ecological feedback, bounded rationality, and calibrated incentives, the interaction among agents may converge toward a stable strategic equilibrium in which environmental preservation becomes individually rational and collectively beneficial.}

\end{itemize}

H5 synthesizes the strategic dimension of the thesis. Preservation becomes stable only if it is compatible with the incentives faced by agents in interaction with one another. The relevant question is therefore not merely whether preservation is socially desirable, but whether it can become a rational strategy in a dynamic setting marked by uncertainty, local ecological constraints, institutional rules, and strategic behavior.

% ---
\subsection{Objectives}
% ---

The objectives of the thesis follow the same structure as the hypotheses. The general objective states the main ambition of the research, while the specific objectives identify the analytical steps required to pursue it.

%...
\subsubsection{General Objective}
%...

The general objective of the thesis is:

\begin{itemize}

    \item[\textbf{G0.}]
    
    \textit{To develop a theoretical, formal, and empirical framework for investigating whether environmental preservation can be incorporated into economic decision-making as an endogenous and incentive-compatible component of production.}

\end{itemize}

This objective is pursued through the articulation of environmental economics, institutional economics, microeconomic theory, dynamic modeling, and empirical analysis applied to Mato Grosso.

%...
\subsubsection{Specific Objectives}
%...

\begin{itemize}

    \item[\textbf{G1.}]
    
    \textit{To examine the theoretical implications of representing the environment as an endogenous cofactor of production rather than merely as an externality, constraint, or fixed stock.}

\end{itemize}

This objective corresponds to H1 and provides the conceptual foundation of the thesis. It seeks to clarify how environmental variables may affect productivity, resilience, opportunity costs, and welfare over time.

\begin{itemize}

    \item[\textbf{G2.}]
    
    \textit{To model environmental decision-making under intertemporal constraints, uncertainty, bounded rationality, incomplete information, and adaptive learning.}

\end{itemize}

This objective corresponds to H2 and supports the formal modeling developed in Chapter 2. It aims to represent the exploitation--preservation dilemma as a dynamic process rather than a static choice.

\begin{itemize}

    \item[\textbf{G3.}]
    
    \textit{To investigate how spatial and ecological scales affect the efficiency and effectiveness of environmental compensation and incentive mechanisms.}

\end{itemize}

This objective corresponds to H3 and examines whether environmental governance should be designed according to broad political boundaries or ecological units such as biomes, watersheds, ecoregions, and transition zones.

\begin{itemize}

    \item[\textbf{G4.}]
    
    \textit{To assess whether positive or mixed incentives can alter the private payoff associated with preservation and reduce reliance on exclusively coercive environmental policy.}

\end{itemize}

This objective corresponds to H4 and evaluates the role of instruments such as carbon credits, payments for environmental services, environmental certification, fiscal incentives, subsidized credit, enforcement, and conditional production rights.

\begin{itemize}

    \item[\textbf{G5.}]
    
    \textit{To analyze whether incentive-compatible preservation can generate stable strategic outcomes under ecological feedback and interaction among heterogeneous agents.}

\end{itemize}

This objective corresponds to H5 and examines whether preservation can become a rational strategy in equilibrium, rather than merely a socially desirable outcome imposed from outside the agent's decision structure.

% ---
\subsection{Methodological Strategy}
% ---

The thesis adopts a mixed research strategy composed of three interdependent methodological layers: theoretical-conceptual analysis, formal modeling, and empirical application. Each layer corresponds to one chapter of the thesis and contributes to a different dimension of the general hypothesis.

The first layer is theoretical and conceptual. It examines the foundations of environmental and resource economics, common-pool resource theory, institutional economics, and production theory in order to identify how the environment has been represented in economic reasoning. This layer introduces the concept of economic commensalism and distinguishes it from predation, extraction, extractivism, and strict preservation. Its purpose is to clarify the conceptual problem that motivates the thesis: the treatment of environmental preservation as an external obligation rather than as a potentially endogenous component of production.

The second layer is formal and deductive. It develops a mathematical model of economic commensalism in which agents make intertemporal decisions under uncertainty, bounded rationality, ecological feedback, and strategic interaction. The model is expected to use dynamic optimization, elements of game theory, stochastic processes, Bayesian updating, Markovian transitions, and filtering mechanisms where appropriate. The purpose of this layer is to examine the internal coherence of the proposed framework and to identify the conditions under which a commensal equilibrium may exist.

The third layer is empirical and applied. It evaluates the plausibility of the theoretical and formal propositions using data from Mato Grosso. The empirical strategy will combine economic, environmental, geospatial, and institutional data to examine the relationship between land use, agricultural productivity, deforestation, ecological heterogeneity, environmental policy, and incentive mechanisms. The precise empirical specification will depend on data availability, policy relevance, and the feasibility of identifying credible counterfactuals.

This methodological structure allows the thesis to proceed from conceptual definition to formal analysis and then to empirical assessment. The theoretical chapter defines the problem and the framework. The mathematical chapter tests the internal consistency of the proposed mechanism. The empirical chapter evaluates whether the mechanism is compatible with observed economic and environmental dynamics in a concrete territory.

% ---
\subsection{Chapter Structure}
% ---

%...
\subsubsection{Chapter 1: Theoretical Framework}
%...

\begin{itemize}

    \item \textbf{Preliminary title:} \emph{Economic Commensalism: A Theoretical Framework to Endogenize Nature in Production Theory}

    \item \textbf{Purpose:} Chapter 1 develops the conceptual foundations of economic commensalism. It reviews the treatment of nature in production theory, environmental economics, common-pool resource theory, and institutional analysis. It also introduces a typology of economy--environment interactions, distinguishing predation, extraction, extractivism, strict preservation, and economic commensalism.

    \item \textbf{Methodological approach:} The chapter adopts a theoretical-analytical methodology based on literature review, conceptual reconstruction, and heuristic modeling. It does not aim to provide empirical proof at this stage, but to establish the conceptual conditions under which environmental preservation can be interpreted as an endogenous and strategically relevant component of production decisions. MUDAR: DEIXAR CLARO QUE GAP EPISTEMOLÓGICO E COERÊNCIA INTERNA VIA FORMALIZAÇÃO SÃO A META.

    \item \textbf{Expected contribution:} The expected contribution of Chapter 1 is to clarify the theoretical gap that motivates the thesis and to provide a coherent conceptual vocabulary for the formal and empirical chapters that follow.

    \item \textbf{Preliminary outline:}

    \begin{itemize}

        \item[i.] \underline{Production theory and the environment}: land, natural resources, externalities, common-pool resources, and the exogenous treatment of ecological systems.

        \item[ii.] \underline{Environmental economics and institutional responses}: Pigouvian taxation, Coasian bargaining, regulation, enforcement, carbon markets, payments for environmental services, certifications, and common-pool governance.

        \item[iii.] \underline{Typology of economy--environment interactions}: predation, extraction, extractivism, strict preservation, and economic commensalism.

        \item[iv.] \underline{Endogenizing environmental variables}: theoretical justification for treating ecological systems as cofactors of production and as determinants of intertemporal productivity.

        \item[v.] \underline{Economic commensalism}: definition, assumptions, theoretical contribution, institutional implications, and limits.

    \end{itemize}

\end{itemize}

%...
\subsubsection{Chapter 2: Formal Model}
%...

\begin{itemize}

    \item \textbf{Preliminary title:} \emph{Economic Commensalism: A Dynamic Production Model with Environmental Endogeneity}

    \item \textbf{Purpose:} Chapter 2 examines the internal consistency of the economic commensalism framework through formal modeling. It investigates whether environmental preservation can become individually rational and strategically stable when ecological variables are incorporated into the agent's intertemporal decision problem. MUDAR: NÃO APENAS EXAMINA, COMO APROFUNDA.

    \item \textbf{Methodological approach:} The chapter adopts a formal-deductive methodology. It is expected to combine dynamic programming, game theory, stochastic processes, Bayesian learning, Markovian dynamics, and filtering mechanisms where these tools are mathematically justified by the structure of the model. The model will identify the conditions under which a commensal equilibrium may exist and compare it with predatory or extractivist baselines.

    \item \textbf{Expected contribution:} The expected contribution of Chapter 2 is to provide a formal representation of the thesis mechanism and to clarify the mathematical conditions required for preservation to be compatible with private rationality and strategic equilibrium.

    \item \textbf{Preliminary outline:}

    \begin{itemize}

        \item[i.] \underline{Basic structure of the model}: agents, state variables, control variables, environmental stock or quality, production, payoff, and constraints.

        \item[ii.] \underline{Intertemporal decision problem}: Bellman equation, value function, preservation choice, exploitation choice, and incentive regimes.

        \item[iii.] \underline{Uncertainty and learning}: ecological shocks, incomplete information, Markovian transitions, Bayesian updating, and informational noise.

        \item[iv.] \underline{Strategic interaction}: interaction among agents, local ecological feedback, governance rules, best responses, and equilibrium conditions.

        \item[v.] \underline{Existence and stability}: continuity, compactness, convexity or concavity conditions, fixed point arguments, convergence, and comparative analysis of incentive regimes.

        \item[vi.] \underline{Computational simulation}: simulation of positive incentives, negative incentives, mixed regimes, local compensation, global compensation, and commensal outcomes compared with extractivist baselines.

    \end{itemize}

\end{itemize}

%...
\subsubsection{Chapter 3: Empirical Application}
%...

\begin{itemize}

    \item \textbf{Preliminary title:} \emph{Economic Efficiency and Environmental Resilience: Empirical Evidence of Commensal Strategies in Mato Grosso}

    \item \textbf{Purpose:} Chapter 3 evaluates the empirical plausibility of the economic commensalism framework using Mato Grosso as a case study. The chapter examines whether economic productivity, land use, deforestation, ecological heterogeneity, and environmental incentives interact in ways consistent with the theoretical and formal propositions developed in Chapters 1 and 2.

    \item \textbf{Methodological approach:} The chapter adopts an empirical strategy combining econometric analysis, geospatial analysis, and institutional interpretation. The precise methods will be selected according to data availability, treatment definition, policy relevance, and identification feasibility. Potential methods include panel data models, difference-in-differences, propensity score matching, synthetic control, spatial econometric models, and geographically weighted regression.

    \item \textbf{Expected contribution:} The expected contribution of Chapter 3 is to assess whether the mechanisms proposed by the thesis are empirically plausible in a territory marked by agricultural expansion, deforestation pressure, ecological heterogeneity, and relevant environmental governance instruments.

    \item \textbf{Preliminary outline:}

    \begin{itemize}

        \item[i.] \underline{Mato Grosso as empirical context}: agricultural expansion, land use, deforestation, biome heterogeneity, institutional arrangements, and environmental policy.

        \item[ii.] \underline{Relevant programs and mechanisms}: carbon markets, REDD+, REDD Early Movers, Environmental Reserve Quotas, payments for environmental services, the Produce, Conserve and Include Program, and other state or federal initiatives.

        \item[iii.] \underline{Data}: agricultural productivity, land use, deforestation, vegetation cover, carbon-related indicators, environmental compliance, rural properties, credit, market access, and socioeconomic controls.

        \item[iv.] \underline{Empirical strategy}: model selection, identification assumptions, spatial scale, temporal scale, treatment definition, control groups, and robustness checks.

        \item[v.] \underline{Expected empirical contribution}: evaluation of whether locally grounded and incentive-compatible mechanisms are associated with better environmental and economic outcomes.

    \end{itemize}

\end{itemize}

% ---
\subsection{Scope and Limitations}
% ---

The thesis is primarily theoretical and exploratory, with a formal and empirical component designed to test the plausibility of the proposed framework. It does not claim that economic commensalism is universally applicable, nor that environmental preservation will always be compatible with private profitability. The argument is conditional: compatibility depends on ecological characteristics, institutional capacity, monitoring technology, incentive design, enforcement credibility, market structure, and the intertemporal preferences of agents.

The empirical chapter is also subject to important limitations. Data availability may restrict the level of spatial disaggregation, the time period analyzed, and the possibility of identifying causal effects. Some mechanisms relevant to economic commensalism, such as expectations, perceived opportunity costs, informal governance, and strategic behavior, may not be directly observable in available datasets. For this reason, the empirical analysis will be interpreted as evidence of plausibility rather than definitive proof of the entire theoretical framework.

A further limitation concerns the normative implications of the thesis. The first three chapters adopt a primarily positive approach: they seek to define, model, and empirically assess the conditions under which preservation may become economically rational. A fourth chapter may be developed later to discuss legal, institutional, and policy applications of the framework, but such normative proposals will depend on the theoretical, formal, and empirical results obtained in the preceding chapters.

% ---
\subsection{Expected Contribution}
% ---

The expected contribution of this thesis is threefold. First, it contributes to environmental and resource economics by proposing a conceptual framework in which ecological preservation is analyzed as a potentially endogenous component of production decisions. Second, it contributes to institutional economics and public policy design by emphasizing the importance of incentive compatibility, local ecological specificity, and intertemporal coordination. Third, it contributes empirically by applying the framework to Mato Grosso, a case in which agricultural productivity, deforestation pressure, ecological heterogeneity, and environmental governance interact in a concrete and policy-relevant setting.

The thesis therefore seeks to advance the discussion on environmental governance by shifting the central question from whether private agents should preserve nature to whether institutions can be designed so that preservation becomes a rational and strategically stable component of economic behavior.

\newpage

% fim %

